\documentclass{article}
\usepackage[T1]{fontenc}
\usepackage{lmodern}
\usepackage{graphicx}
\usepackage{amsmath,amssymb}
\usepackage[a4paper,margin=2cm]{geometry}
\usepackage[absolute,overlay]{textpos}
\setlength{\TPHorizModule}{1mm}
\setlength{\TPVertModule}{1mm}
\usepackage[indonesian]{babel}
\usepackage{hyperref}
\setlength{\emergencystretch}{2em}
% unit: Halaman 34

\begin{document}

% === Halaman 34 (llm) ===

(b) $2P = R$

\[ m \equiv \frac{3x_p^2 + a}{2y_p} \pmod{p} \]

\begin{align*}
 m &= (3(2)^2 + 1)/8 \pmod{11} \equiv 13/8 \pmod{11} \\
 &\equiv 13 \cdot 8^{-1} \pmod{11} \\
 &\equiv 13 \cdot 7 \pmod{11} \\
 &\equiv 78 \pmod{11} \equiv 3 \pmod{11} 
\end{align*}

Koordinat R:

\begin{align*}
 x_r &\equiv m^2 - 2x_p \pmod{p} \equiv 3^2 - 2 \cdot 2 \pmod{11} \equiv 5 \pmod{11} \\
 y_r &\equiv m(x_p - x_r) - y_p \pmod{p} \equiv 3(2 - 5) - 4 \pmod{11} \\
 &\equiv -13 \pmod{11} \equiv 9 \pmod{11} 
\end{align*}

Jadi, $R(5, 9)$

\end{document}
